\chapter{2014年广东卷（理）}

\section{单选题}

\begin{problem}
已知集合 \(M = \{-1, 0, 1\}\)，\(N = \{0, 1, 2\}\)，则 \(M \cup N =\)（\quad）

\choices
    {\(\{0, 1\}\)}
    {\(\{-1, 0, 2\}\)}
    {\(\{-1, 0, 1, 2\}\)}
    {\(\{-1, 0, 1\}\)}
\end{problem}

\begin{answer}
C
\end{answer}

\begin{solution}
并集包含两个集合中的全部元素，因此
\(M\cup N=\{-1,0,1,2\}\)，与选项 C 相符，故选 C.
\end{solution}

\begin{problem}
已知复数 \(z\) 满足 \((3 + 4\i)z = 25\)，则 \(z =\)（\quad）

\choices
    {\(-3 + 4\i\)}
    {\(-3 - 4\i\)}
    {\(3 + 4\i\)}
    {\(3 - 4\i\)}
\end{problem}

\begin{answer}
D
\end{answer}

\begin{solution}
将分母实数化，得
\(z=\dfrac{25}{3+4\i}=\dfrac{25(3-4\i)}{(3+4\i)(3-4\i)}=3-4\i\)，与选项 D 相符，故选 D.
\end{solution}

\begin{problem}
若变量 \(x\)，\(y\) 满足约束条件 \(\begin{cases}
y \leqslant x,\\
x + y \leqslant 1,\\
y \geqslant -1,
\end{cases}\) 且 \(z = 2x + y\) 的最大值和最小值分别为 \(m\) 和 \(n\)，则 \(m - n =\)（\quad）

\choices
    {\(5\)}
    {\(6\)}
    {\(7\)}
    {\(8\)}
\end{problem}

\begin{answer}
B
\end{answer}

\begin{solution}
三个边界的交点分别为
\((-1,-1)\)、\((2,-1)\) 和 \(\left(\frac{1}{2},\frac{1}{2}\right)\)，可行域是由这三个点围成的三角形. 因为线性目标函数在有界多边形上的最值在顶点处取得，所以
\(2(-1)+(-1)=-3\)，\(2\cdot2+(-1)=3\)，\(2\cdot\frac{1}{2}+\frac{1}{2}=\frac{3}{2}\).
故最大值 \(m=3\)，最小值 \(n=-3\)，从而 \(m-n=6\)，故选 B.
\end{solution}

\begin{problem}
若实数 \(k\) 满足 \(0 < k < 9\)，则曲线 \(\frac{x^{2}}{25} - \frac{y^{2}}{9-k} = 1\) 与曲线 \(\frac{x^{2}}{25-k} - \frac{y^{2}}{9} = 1\) 的（\quad）

\choices
    {焦距相等}
    {实半轴长相等}
    {虚半轴长相等}
    {离心率相等}
\end{problem}

\begin{answer}
A
\end{answer}

\begin{solution}
第一条双曲线的实半轴长和虚半轴长分别为 \(5\) 和 \(\sqrt{9-k}\)，所以其焦半距满足
\(c_1^2=5^2+(\sqrt{9-k})^2=34-k\). 第二条双曲线的实半轴长和虚半轴长分别为 \(\sqrt{25-k}\) 和 \(3\)，所以
\(c_2^2=(\sqrt{25-k})^2+3^2=34-k\). 因此 \(c_1=c_2\)，两条双曲线的焦距 \(2c\) 相等，故选 A.
\end{solution}

\begin{problem}
已知向量 \(\bs{a} = (1, 0, -1)\)，则下列向量中与 \(\bs{a}\) 成 \(60^{\circ}\) 角的是（\quad）

\choices
    {\((-1, 1, 0)\)}
    {\((1, -1, 0)\)}
    {\((0, -1, 1)\)}
    {\((-1, 0, 1)\)}
\end{problem}

\begin{answer}
B
\end{answer}

\begin{solution}
\(\lvert\bs{a}\rvert=\sqrt{1^2+0^2+(-1)^2}=\sqrt{2}\). 逐一计算各选项向量 \(v\) 与 \(\bs{a}\) 的夹角余弦：
\(\frac{\bs{a}\cdot(-1,1,0)}{|\bs{a}|\sqrt{2}}=-\frac12\)、\(\frac{\bs{a}\cdot(1,-1,0)}{|\bs{a}|\sqrt{2}}=\frac12\)、\(\frac{\bs{a}\cdot(0,-1,1)}{|\bs{a}|\sqrt{2}}=-\frac12\)、\(\frac{\bs{a}\cdot(-1,0,1)}{|\bs{a}|\sqrt{2}}=-1\).
只有第二个向量的夹角余弦为 \(\cos60^\circ=\frac12\)，故选 B.
\end{solution}

\begin{problem}
已知某地区中小学生人数和近视情况如图 1 和图 2 所示，为了解该地区中小学生的近视形成原因，用分层抽样的方法抽取 \(2\%\) 的学生进行调查，则样本容量和抽取的高中生近视人数分别为（\quad）

\bitmapfigure[max width=0.85\linewidth]{img/2014/guangdong_science/q6_fig1.png}

\choices
    {\(200,20\)}
    {\(100,20\)}
    {\(200,10\)}
    {\(100,10\)}
\end{problem}

\begin{answer}
A
\end{answer}

\begin{solution}
该地区学生总人数为 \(3500+4500+2000=10000\)，抽取 \(2\%\) 后样本容量为
\(10000\times2\%=200\). 图 2 表明高中生近视率为 \(50\%\)，高中生近视人数为 \(2000\times50\%=1000\)，按同样比例抽取的高中生近视人数为 \(1000\times2\%=20\). 因此选 A.
\end{solution}

\begin{problem}
若空间中四条两两不同的直线 \(l_{1}\)，\(l_{2}\)，\(l_{3}\)，\(l_{4}\)，满足 \(l_{1} \perp l_{2}\)，\(l_{2} \perp l_{3}\)，\(l_{3} \perp l_{4}\)，则下列结论一定正确的是（\quad）

\choices
    {\(l_{1} \perp l_{4}\)}
    {\(l_{1} \parallel l_{4}\)}
    {\(l_{1}\)，\(l_{4}\) 既不垂直也不平行}
    {\(l_{1}\)，\(l_{4}\) 的位置关系不确定}
\end{problem}

\begin{answer}
D
\end{answer}

\begin{solution}
例如取三条相交坐标轴所在的直线
\(l_2\) 为 \(x\) 轴，\(l_1\) 为 \(y\) 轴，\(l_3\) 为 \(z\) 轴. 若取过 \((0,0,1)\) 且平行于 \(y\) 轴的 \(l_4\)，则 \(l_1\parallel l_4\)，且题设的三个垂直关系都成立. 若改取过原点、方向向量为 \((1,1,0)\) 的 \(l_4\)，则 \(l_3\perp l_4\)，但 \(l_1\) 与 \(l_4\) 相交且不垂直，也不平行. 故 \(l_1\) 与 \(l_4\) 的位置关系不确定，选 D.
\end{solution}

\begin{problem}
设集合 \(A = \{(x_{1}, x_{2}, x_{3}, x_{4}, x_{5}) \mid x_{i} \in \{-1, 0, 1\}, i = 1, 2, 3, 4, 5\}\)，那么集合 \(A\) 中满足条件“\(1 \leqslant |x_{1}| + |x_{2}| + |x_{3}| + |x_{4}| + |x_{5}| \leqslant 3\)”的元素个数为（\quad）

\choices
    {\(60\)}
    {\(90\)}
    {\(120\)}
    {\(130\)}
\end{problem}

\begin{answer}
D
\end{answer}

\begin{solution}
令恰有 \(r\) 个分量非零. 因为每个非零分量只能取 \(1\) 或 \(-1\)，所以当 \(r=1\)，\(2\)，\(3\) 时，对应元素个数分别为
\(\mathrm C_5^1 2^1=10\)、\(\mathrm C_5^2 2^2=40\)、\(\mathrm C_5^3 2^3=80\). 题设条件正是 \(r=1\)，\(2\)，\(3\)，因此元素总数为
\(10+40+80=130\)，故选 D.
\end{solution}

\section{填空题}

\begin{problem}
不等式 \( |x - 1| + |x + 2| \geqslant 5 \) 的解集为\fillinblank{}.
\end{problem}

\begin{answer}
\(\left(-\infty,-3\right]\cup\left[2,+\infty\right)\)
\end{answer}

\begin{solution}
按绝对值的分界点 \(-2\)、\(1\) 分类. 当 \(x\leqslant-2\) 时，原不等式化为
\(-2x-1\geqslant5\)，得 \(x\leqslant-3\)；当 \(-2\leqslant x\leqslant1\) 时，左端恒为 \(3\)，无解；当 \(x\geqslant1\) 时，原不等式化为 \(2x+1\geqslant5\)，得 \(x\geqslant2\). 由于原不等式是非严格不等式，端点应保留，所以解集为
\(\left(-\infty,-3\right]\cup\left[2,+\infty\right)\).
\end{solution}

\begin{problem}
曲线 \( y = \e^{-5x} + 2 \) 在点 \((0, 3)\) 处的切线方程为\fillinblank{}.
\end{problem}

\begin{answer}
\(y=-5x+3\)
\end{answer}

\begin{solution}
设 \(g(x)=\e^{-5x}+2\)，则 \(g'(x)=-5\e^{-5x}\)，所以 \(g'(0)=-5\). 切线经过 \((0,3)\)，故其方程为
\(y-3=-5(x-0)\)，即 \(y=-5x+3\).
\end{solution}

\begin{problem}
从 0， 1， 2， 3， 4， 5， 6， 7， 8， 9 中任取七个不同的数，则这七个数的中位数是 6 的概率为\fillinblank{}.
\end{problem}

\begin{answer}
\(\frac{1}{6}\)
\end{answer}

\begin{solution}
七个数的中位数是第 4 个数. 要使中位数为 \(6\)，必须选取数字 \(6\)，再从 \(0\)，\(1\)，\(2\)，\(3\)，\(4\)，\(5\) 中选出三个，从 \(7\)，\(8\)，\(9\) 中选出三个. 因此有利选法数为
\(\mathrm C_6^3\mathrm C_3^3=20\). 全部选法数为 \(\mathrm C_{10}^7=120\)，所以所求概率为
\(\dfrac{20}{120}=\dfrac16\).
\end{solution}

\begin{problem}
在 \(\triangle ABC\) 中，角 \(A\)，\(B\)，\(C\) 所对应的边分别为 \(a\)，\(b\)，\(c\)，已知 \(b \cos C + c \cos B = 2b\)，则 \(\frac{a}{b} =\)\fillinblank{}.
\end{problem}

\begin{answer}
2
\end{answer}

\begin{solution}
由余弦定理，
\(b\cos C=\dfrac{a^2+b^2-c^2}{2a}\)，\(c\cos B=\dfrac{a^2+c^2-b^2}{2a}\). 相加得
\[
b\cos C+c\cos B=\frac{2a^2}{2a}=a
\]
题设给出该和等于 \(2b\)，所以 \(a=2b\)，从而 \(\dfrac ab=2\).
\end{solution}

\begin{problem}
若等比数列 \(\{a_{n}\}\) 的各项均为正数，且 \(a_{10}a_{11} + a_{9}a_{12} = 2\e^{5}\)，则 \(\ln a_1 + \ln a_2 + \cdots + \ln a_{20} =\fillinblank{}\).
\end{problem}

\begin{answer}
50
\end{answer}

\begin{solution}
设等比数列的首项为 \(a_1\)，公比为 \(q\)，则 \(q>0\)，且
\(a_{10}a_{11}=a_1^2q^{19}=a_9a_{12}\). 由题设
\(2a_1^2q^{19}=2\e^5\)，即 \(a_1^2q^{19}=\e^5\). 于是
\[
\ln a_1+\cdots+\ln a_{20}
=\ln\left(a_1a_2\cdots a_{20}\right)
=\ln\left(a_1^{20}q^{0+1+\cdots+19}\right)
=\ln\left((a_1^2q^{19})^{10}\right)=\ln\e^{50}=50
\]
\end{solution}

\section{选考题：填空题}

\begin{problem}
在极坐标系中，曲线 \(C_1\) 和 \(C_2\) 的方程分别为 \(\rho \sin^2 \theta = \cos \theta\) 和 \(\rho \sin \theta = 1\)，以极点为平面直角坐标系的原点，极轴为 \(x\) 轴的正半轴，建立平面直角坐标系，则曲线 \(C_1\) 和 \(C_2\) 的交点的直角坐标为\fillinblank{}.
\end{problem}

\begin{answer}
\((1,1)\)
\end{answer}

\begin{solution}
由极坐标与直角坐标的关系 \(x=\rho\cos\theta\)、\(y=\rho\sin\theta\)，第一条曲线满足
\(\rho^2\sin^2\theta=\rho\cos\theta\)，即 \(y^2=x\)；第二条曲线满足 \(y=1\). 代入得 \(x=1\)，所以交点的直角坐标为 \((1,1)\).
\end{solution}

\begin{problem}
如图，在平行四边形 \(ABCD\) 中，点 \(E\) 在 \(AB\) 上，且 \(EB = 2AE\)，\(AC\) 与 \(DE\) 交于点 \(F\)，则 \(\triangle CDF\) 的面积与 \(\triangle AEF\) 的面积之比为\fillinblank{}.

\bitmapfigure[max width=0.45\linewidth]{img/2014/guangdong_science/q15_fig1.png}
\end{problem}

\begin{answer}
9
\end{answer}

\begin{solution}
取 \(AE=1\)、\(EB=2\)，建立直角坐标系，使 \(A=(0,0)\)、\(B=(3,0)\)、\(D=(u,v)\)（\(v>0\)），则 \(C=(3+u,v)\)、\(E=(1,0)\). 设 \(F=AC\cap DE\)，在两条直线上分别写成
\(F=t(3+u,v)\) 和 \(F=(1+s(u-1),sv)\). 比较纵坐标得 \(s=t\)，再比较横坐标得
\(t(3+u)=1+t(u-1)\)，从而 \(t=\frac14\).

因此 \(F\) 到 \(AB\) 的距离是 \(\frac14\) 个平行四边形高，而 \(C\)、\(D\) 到直线 \(AB\) 的距离是一个平行四边形高. 于是 \(S_{\triangle AEF}=\frac12\cdot AE\cdot\frac{h}{4}=\frac h8\)，\(S_{\triangle CDF}=\frac12\cdot CD\cdot\frac{3h}{4}=\frac{9h}{8}\).
其中 \(CD=AB=3AE\). 故
\(\dfrac{S_{\triangle CDF}}{S_{\triangle AEF}}=9\).
\end{solution}

\section{解答题}

\begin{problem}
已知函数 \(f(x) = A \sin \left( x + \frac{\pi}{4} \right)\)，\(x \in \R\)，且 \( f\left( \frac{5\pi}{12} \right) = \frac{3}{2} \).

\begin{enumerate}
\item 求 \(A\) 的值；
\item 若 \(f(\theta) + f(-\theta) = \frac{3}{2}\)，\(\theta \in \left(0, \frac{\pi}{2}\right)\)，求 \( f\left(\frac{3}{4}\pi - \theta\right) \).
\end{enumerate}
\end{problem}

\begin{solution}
\begin{enumerate}
\item 由题设
\[
f\left(\frac{5\pi}{12}\right)=A\sin\left(\frac{5\pi}{12}+\frac{\pi}{4}\right)=A\sin\frac{2\pi}{3}=\frac{\sqrt{3}}{2}A=\frac32
\]
所以 \(A=\sqrt3\)，从而 \(f(x)=\sqrt3\sin\left(x+\frac\pi4\right)\).

\item 由和差化积公式，
\[
f(\theta)+f(-\theta)
=\sqrt3\left[\sin\left(\theta+\frac\pi4\right)+\sin\left(\frac\pi4-\theta\right)\right]
=\sqrt6\cos\theta=\frac32
\]
故 \(\cos\theta=\frac{\sqrt6}{4}\). 因为 \(0<\theta<\frac\pi2\)，所以
\(\sin\theta=\sqrt{1-\frac{6}{16}}=\frac{\sqrt{10}}{4}\). 最后
\[
f\left(\frac{3\pi}{4}-\theta\right)
=\sqrt3\sin(\pi-\theta)
=\sqrt3\sin\theta
=\frac{\sqrt{30}}{4}
\]
\end{enumerate}
\end{solution}

\begin{problem}
随机观测生产某种零件的某工厂 25 名工人的日加工零件数（单位：件），获得数据如下：30，42，41，36，44，40，37，37，25，45，29，43，31，36，49，34，33，43，38，42，32，34，46，39，36.
根据上述数据得到样本的频率分布表如下：

\begin{center}
\begin{tabular}{c|cc}
\hline
分组 & 频数 & 频率 \\
\hline
\([25,30]\) & 3 & 0.12 \\
\((30,35]\) & 5 & 0.20 \\
\((35,40]\) & 8 & 0.32 \\
\((40,45]\) & \(n_1\) & \(f_1\) \\
\((45,50]\) & \(n_2\) & \(f_2\) \\
\hline
\end{tabular}
\end{center}

\begin{enumerate}
\item 确定样本频率分布表中 \(n_1\)，\(n_2\)，\(f_1\) 和 \(f_2\) 的值；
\item 根据上述频率分布表，画出样本频率分布直方图；
\item 根据样本频率分布直方图，求在该厂任取 4 人，至少有 1 人的日加工零件数落在区间 \((30,35]\) 的概率.
\end{enumerate}
\end{problem}

\begin{solution}
\begin{enumerate}
\item 直接统计原始数据，各组频数分别为 \(3\)，\(5\)，\(8\)，\(7\)，\(2\)，所以
\(n_1=7\)、\(n_2=2\)，并且
\(f_1=\frac{7}{25}=0.28\)、\(f_2=\frac{2}{25}=0.08\).

 \item 直方图的纵轴表示频率除以组距. 五组组距均为 \(5\)，所以从左到右五个矩形的高分别为
\(\frac{0.12}{5}=0.024\)、\(\frac{0.20}{5}=0.04\)、\(\frac{0.32}{5}=0.064\)、\(\frac{0.28}{5}=0.056\)、\(\frac{0.08}{5}=0.016\).
因此以 \([25,30]\)、\((30,35]\)、\((35,40]\)、\((40,45]\)、\((45,50]\) 为底边，以上述五个数为高，分别作相邻矩形，即得样本频率分布直方图.

 \item 由直方图，任意抽取一名工人的日加工零件数落在 \((30,35]\) 的概率用其频率估计为 \(0.20\)，不落在该区间的概率为 \(0.80\). 四人中至少有一人落在该区间的概率为
\(1-0.80^4=1-0.4096=0.5904\).
\end{enumerate}
\end{solution}

\begin{problem}
如图，四边形 \(ABCD\) 为正方形，\(PD \perp\) 平面 \(ABCD\)，\(\angle DPC = 30^{\circ}\)，\(AF \perp PC\) 于点 \(F\)，\(FE \parallel CD\)，交 \(PD\) 于点 \(E\).

\begin{enumerate}
\item 证明：\(CF \perp\) 平面 \(ADF\)；
\item 求二面角 \(D - AF - E\) 的余弦值.
\end{enumerate}

\bitmapfigure[max width=0.56\linewidth]{img/2014/guangdong_science/q18_fig1.png}
\end{problem}

\begin{solution}
\begin{enumerate}
\item
取正方形边长 \(DC=1\)，建立空间直角坐标系，使
\(D=(0,0,0)\)、\(C=(1,0,0)\)、\(A=(0,1,0)\)、\(B=(1,1,0)\). 因为 \(PD\perp\) 平面 \(ABCD\)，且三角形 \(PDC\) 在 \(D\) 处为直角，故
\(\tan\angle DPC=\frac{DC}{PD}=\tan30^\circ\)，从而 \(P=(0,0,\sqrt3)\).

点 \(F\) 在 \(PC\) 上，设 \(F=(t,0,\sqrt3(1-t))\). 由 \(AF\perp PC\)，有
\[
\overrightarrow{AF}\cdot\overrightarrow{PC}
=\left(t,-1,\sqrt3(1-t)\right)\cdot(1,0,-\sqrt3)
=4t-3=0
\]
所以 \(F=\left(\frac34,0,\frac{\sqrt3}{4}\right)\). 又因 \(FE\parallel CD\)，且 \(E\in PD\)，故 \(E=\left(0,0,\frac{\sqrt3}{4}\right)\).

计算得
\(\overrightarrow{CF}=\left(-\frac14,0,\frac{\sqrt3}{4}\right)\)、
\(\overrightarrow{DF}=\left(\frac34,0,\frac{\sqrt3}{4}\right)\)、
\(\overrightarrow{AF}=\left(\frac34,-1,\frac{\sqrt3}{4}\right)\)，于是
\(\overrightarrow{CF}\cdot\overrightarrow{DF}=0\)，且 \(\overrightarrow{CF}\cdot\overrightarrow{AF}=0\). 直线 \(AF\)、\(DF\) 是平面 \(ADF\) 内相交的两条直线，所以 \(CF\perp\) 平面 \(ADF\).

\item 由 \(\overrightarrow{AD}=(0,-1,0)\)、\(\overrightarrow{AF}=\left(\frac34,-1,\frac{\sqrt3}{4}\right)\)，以及 \(\overrightarrow{AE}=\left(0,-1,\frac{\sqrt3}{4}\right)\)，可分别取平面 \(ADF\) 和平面 \(EAF\) 的法向量为
\(\bs n_1=(\sqrt3,0,-3)\) 和 \(\bs n_2=(0,\sqrt3,4)\). 设二面角 \(D-AF-E\) 为 \(\varphi\)，则取对应的锐角余弦得
\[
\cos\varphi=\frac{|\bs n_1\cdot\bs n_2|}{|\bs n_1||\bs n_2|}
=\frac{12}{\sqrt{12}\sqrt{19}}
=\frac{2\sqrt{57}}{19}
\]
\end{enumerate}
\end{solution}

\begin{problem}
设数列 \(\{a_{n}\}\) 的前 \(n\) 和为 \(S_{n}\)，满足 \(S_{n} = 2na_{n+1} - 3n^{2} - 4n\)，\(n \in \N^*\)，且 \(S_{3} = 15\).

\begin{enumerate}
\item 求 \( a_{1} \)，\( a_{2} \)，\( a_{3} \) 的值；
\item 求数列 \(\{a_{n}\}\) 的通项公式.
\end{enumerate}
\end{problem}

\begin{solution}
\begin{enumerate}
\item
在 \(n=2\) 时，题设给出 \(S_2=4a_3-20\). 又
\(S_3=S_2+a_3=15\)，所以 \(5a_3-20=15\)，得 \(a_3=7\)，进而 \(S_2=8\).

在 \(n=1\) 时，\(S_1=2a_2-7\)，而 \(S_2=S_1+a_2=8\)，所以 \(3a_2-7=8\)，得 \(a_2=5\)，再得 \(a_1=S_1=3\).

\item 下面用数学归纳法求通项. 已知 \(a_1=3=2\cdot1+1\). 假设对某个正整数 \(n\)，前 \(n\) 项均满足 \(a_j=2j+1\)，则
\(S_n=3+5+\cdots+(2n+1)=n(n+2)\). 代入已知关系，得
\[
n(n+2)=2na_{n+1}-3n^2-4n
\]
由于 \(n\geqslant1\)，除以 \(n\) 后得
\(n+2=2a_{n+1}-3n-4\)，即 \(a_{n+1}=2n+3=2(n+1)+1\). 因此对任意正整数 \(n\)，都有
\(a_n=2n+1\).
\end{enumerate}
\end{solution}

\begin{problem}
已知椭圆 \(C: \frac{x^2}{a^2} + \frac{y^2}{b^2} = 1\)（\(a > b > 0\)）的一个焦点为 \((\sqrt{5}, 0)\). 离心率为 \(\frac{\sqrt{5}}{3}\).

\begin{enumerate}
\item 求椭圆 \(C\) 的标准方程；
\item 若动点 \(P(x_0, y_0)\) 为椭圆 \(C\) 外一点，且点 \(P\) 到椭圆 \(C\) 的两条切线相互垂直，求点 \(P\) 的轨迹方程.
\end{enumerate}
\end{problem}

\begin{solution}
\begin{enumerate}
\item 椭圆焦半距 \(c=\sqrt5\)，又 \(e=\frac ca=\frac{\sqrt5}{3}\)，所以 \(a=3\). 由 \(a^2=b^2+c^2\)，得 \(b^2=9-5=4\)，故椭圆标准方程为
\(\dfrac{x^2}{9}+\dfrac{y^2}{4}=1\).

\item 先设一般直线过 \(P(x_0,y_0)\) 且斜率为 \(m\)，写成 \(y-y_0=m(x-x_0)\). 令 \(q=y_0-mx_0\)，代入椭圆方程得
\(\frac{x^2}{9}+\frac{(mx+q)^2}{4}=1\)，即
\((4+9m^2)x^2+18mqx+9q^2-36=0\). 直线为切线时该关于 \(x\) 的二次方程有重根，故
\((18mq)^2-4(4+9m^2)(9q^2-36)=0\)，化简得 \(q^2=4+9m^2\). 代入 \(q=y_0-mx_0\)，得
\[
(x_0^2-9)m^2-2x_0y_0m+y_0^2-4=0
\]
两条切线的斜率 \(m_1\)，\(m_2\) 是此方程的两个根. 当两条切线均非竖直时，垂直条件为 \(m_1m_2=-1\)，而根的乘积为 \(\frac{y_0^2-4}{x_0^2-9}\)，所以
\[
\frac{y_0^2-4}{x_0^2-9}=-1
\]
从而 \(x_0^2+y_0^2=13\).
若有一条切线为竖直线，则另一条必须为水平线，这对应于 \(P=(\pm3,\pm2)\)，这些点同样满足上式. 反之，圆 \(x_0^2+y_0^2=13\) 上的点满足 \(\frac{x_0^2}{9}+\frac{y_0^2}{4}\geqslant\frac{x_0^2+y_0^2}{9}=\frac{13}{9}>1\)，确为椭圆外点. 因此轨迹方程为
\(x_0^2+y_0^2=13\).
\end{enumerate}
\end{solution}

\begin{problem}
设函数 \( f(x) = \frac{1}{\sqrt{(x^2 + 2x + k)^2 + 2(x^2 + 2x + k) - 3}} \)，其中 \( k < -2 \).

\begin{enumerate}
\item 求函数 \( f(x) \) 的定义域 \(D\)（用区间表示）；
\item 讨论 \(f(x)\) 在区间 \(D\) 上的单调性；
\item 若 \(k < -6\)，求 \(D\) 上满足条件 \(f(x) > f(1)\) 的 \(x\) 的集合（用区间表示）.
\end{enumerate}
\end{problem}

\begin{solution}
\begin{enumerate}
\item
令 \(u=x^2+2x+k\)，则根式内为 \(u^2+2u-3=(u-1)(u+3)\).

函数有意义要求根式内严格大于零，因此 \(u<-3\) 或 \(u>1\). 分别化为
\( (x+1)^2<-2-k\) 或 \((x+1)^2>2-k\)，所以
\[
D=\left(-\infty,-1-\sqrt{2-k}\right)\cup\left(-1-\sqrt{-2-k},-1+\sqrt{-2-k}\right)\cup\left(-1+\sqrt{2-k},+\infty\right)
\]

\item 在定义域内
\[
f'(x)=-\frac{2(x+1)(u+1)}{\left(u^2+2u-3\right)^{3/2}}
\]
在两个外侧区间上 \(u>1\)，所以 \(u+1>0\)，导数符号与 \(-(x+1)\) 相同，故左外侧递增、右外侧递减. 在中间区间上 \(u<-3\)，所以 \(u+1<0\)，导数符号与 \(x+1\) 相同，故
\(f(x)\) 在 \(\left(-1-\sqrt{-2-k},-1\right)\) 上递减，在 \(\left(-1,-1+\sqrt{-2-k}\right)\) 上递增. 综上，递增区间为
\(\left(-\infty,-1-\sqrt{2-k}\right)\) 和 \(\left(-1,-1+\sqrt{-2-k}\right)\)，递减区间为
\(\left(-1-\sqrt{-2-k},-1\right)\) 和 \(\left(-1+\sqrt{2-k},+\infty\right)\).

\item 当 \(k<-6\) 时，\(u(1)=k+3<-3\)，且在定义域内分母为正. 由
\(f(x)>f(1)\) 等价于
\[
(u+1)^2<(u(1)+1)^2=(k+4)^2
\]
因为 \(k<-6\)，所以这等价于 \(k+3<u<-k-5\). 再与定义域中的 \(u<-3\) 或 \(u>1\) 结合，得到
\(k+3<u<-3\) 或 \(1<u<-k-5\).

当 \(k+3<u<-3\) 时，等价于 \(x<-3\) 或 \(x>1\)，且 \((x+1)^2<-2-k\)，得到
\(\left(-1-\sqrt{-2-k},-3\right)\cup\left(1,-1+\sqrt{-2-k}\right)\).

当 \(1<u<-k-5\) 时，等价于 \((x+1)^2>2-k\) 且 \((x+1)^2<-2k-4\)，得到
\(\left(-1-\sqrt{-2k-4},-1-\sqrt{2-k}\right)\cup\left(-1+\sqrt{2-k},-1+\sqrt{-2k-4}\right)\).

所以所求集合为
\[
\left(-1-\sqrt{-2k-4},-1-\sqrt{2-k}\right)\cup\left(-1-\sqrt{-2-k},-3\right)
\]
和
\[
\left(1,-1+\sqrt{-2-k}\right)\cup\left(-1+\sqrt{2-k},-1+\sqrt{-2k-4}\right)
\]
\end{enumerate}
\end{solution}
